\section{Functoriality and further elliptic branches}
\label{rev:sec:elliptic-extensions}

The classification is stable under an algebraic change of variable only
after its branch, derivative, and convergence domain have been specified.
The following elementary observation keeps these three issues separate.
For a germ $F$ and a pullback $G(t)=r(t)F(\phi(t))$, put
\[
 \delta=t\frac{d}{dt},\qquad
 u=\delta\log r,\qquad v=\delta\log\phi.
\]
At a regular point with $rv\ne0$, the first jets satisfy
\begin{equation}\label{rev:eq:jet-pullback}
 \begin{pmatrix}G\\ \delta G\end{pmatrix}
 =r\begin{pmatrix}1&0\\u&v\end{pmatrix}
   \begin{pmatrix}F\circ\phi\\(\thetaop F)\circ\phi\end{pmatrix}.
\end{equation}
Thus pullback induces an isomorphism of coefficient lines,
\[
 [A:B]\longmapsto[A+Bu:Bv].
\]
Algebraicity, the period reduction, and projective uniqueness transport
through this map. An identity of germs transports analytic values;
ordinary convergence of the resulting series is a separate assertion.

\subsection{Quadratic companions}\label{sec:quadratic-companion-extra}
Put $K=R_\ell/4$, $\mathcal F(X)=F_{s_\ell}(R_\ell X)
=\sum u_\ell(n)X^n$, and define
\[
 U(n)=\sum_{k=0}^n u_\ell(k)\binom{n+k}{n-k}(-K)^{n-k},
 \qquad \mathcal G(Y)=\sum_{n\ge0}U(n)Y^n.
\]
The binomial generating function gives
\begin{equation}\label{eq:quad-extra-generating}
 \mathcal G(Y)=\frac1{1+KY}
      \mathcal F\!\left(\frac{Y}{(1+KY)^2}\right).
\end{equation}
For $c\ne0,-K$ set $J=(c+K)^2/c$. Formula
\eqref{rev:eq:jet-pullback} becomes
\begin{equation}\label{eq:quad-extra-transport}
 (A+B\delta)\mathcal G(1/c)
 =\frac1J\bigl((c+K)A-KB+(c-K)B\Theta\bigr)
                         \mathcal F^{\mathrm{an}}(1/J),
 \qquad \Theta=X\partial_X.
\end{equation}
The continuation is along the real interval from zero. For $|c|>K$
the coefficient map is invertible. Its image has $J>4K$ if $c>K$
and $J<0$ if $c<-K$; in the latter case the standard power series
need not converge.

\begin{proposition}[The five quadratic companions]
\label{prop:quad-extra-five}
Among ordinarily convergent quadratic-companion identities with
$c\in\Z\setminus\{0\}$, $\gcd(A,B)=1$, and a CM elliptic fiber,
the following are precisely the classes up to simultaneous sign:
\[
\begin{array}{c|rrr|l}
 \ell&A&B&c&\alpha\\\hline
 2&2&9&512&4\\
 2&0&1&-256&4\sqrt3/9\\
 2&1&9&-4096&64\sqrt7/49\\
 3&1&8&-243&3\sqrt3\\
 4&0&1&-32&2.
\end{array}
\]
Here the value is $\alpha/\pi$, and every displayed sum is absolutely
convergent.
\end{proposition}
\begin{proof}
Clausen's and Pfaff's transformations give
\[
 \mathcal G(Y)=(1+KY)^{2s-1}{}_2F_1(s,s;1;-KY)^2,
 \qquad s=s_\ell.
\]
The only singularity on $|KY|=1$ is $KY=-1$. The Gauss connection
formula there gives, for $s<1/2$,
\[
 \frac{U(n)}{(-K)^n}
 =\frac{\Gamma(1-2s)}{\Gamma(1-s)^4}n^{-2s}
       +O(n^{-1-\epsilon_s}),\qquad \epsilon_s>0;
\]
for $s=1/2$ it gives $(d\log n+e)/n+O(\log n/n^2)$ with $d>0$.
Consequently a nonzero linear weight converges absolutely for $|c|>K$,
diverges for $|c|<K$ or $c=-K$, and at $c=K$ converges exactly when
$B=0$. Abel's theorem identifies the last value with
$\mathcal G(1/K)=F_s(1)/2$, excluded by the positive-endpoint theorem.

At CM, Proposition~\ref{prop:rational-cm-integrality} makes the rational
standard denominator $J$ integral. Since
\[
 J=c+2K+\frac{K^2}{c},
\]
we have $c\mid K^2$. Thus integrality converts the classification into
a finite divisor problem. The signed divisors satisfying $|c|>K$ number
$62,12,6,8$ at the four levels. Applying the class-polynomial criterion
of Section~\ref{sec:enumeration} gives exactly
\[
 (\ell,c,J)=(2,512,648),\ (2,-256,-144),\ (2,-4096,-3969),
             \ (3,-243,-192),\ (4,-32,-8).
\]
Their selected elliptic invariants are $1728,0,-3375,0,1728$,
respectively. This is the same exact finite test as for the standard
tables, applied before a standard convergence filter.

Three evaluations follow from \eqref{eq:quad-extra-transport}: the
transported weights are respectively $8/9$, $64/441$, and $9/4$ times
the standard pairs $(1,7)$, $(8,65)$, and $(1,5)$, at
$J=648,-3969,-192$. The two remaining evaluations follow from the
endomorphisms of the equianharmonic and square lattices. The period
normalization and determinant identity continue along the nonsingular
real interval $z<1$, so their use does not require convergence of the
transported standard power series. The complete data, in the notation
of \eqref{eq:cm-matrix-alpha}, are
\[
\begin{array}{c|c|c|c|c|c}
 (\ell,c)&(a,b)&(u,v)&\lambda&\kappa&d\\\hline
 (2,-256)&(4/9,20/9)&(-1/5,-1/5)&\zeta&0&1\\
 (4,-32)&(2,6)&(-1/3,-2/9)&i&0&1
\end{array}
 \qquad \zeta=\frac{-1+i\sqrt3}{2}.
\]
Indeed the fibers are $y^2=4X^3-4$ and $y^2=4X^3-9X$, with
endomorphisms $(X,y)\mapsto(\zeta X,y)$ and
$(X,y)\mapsto(-X,iy)$. Their standard cycles are the positive primitive
real cycles, transported from the primitive vanishing cycle; the
oriented lattice bases are $(\omega,\zeta\omega)$ and
$(\omega,i\omega)$. Since $a+bu=0$, the determinant formula yields
$4\sqrt3/9$ and $2$.

Finally, the period-ratio proof of uniqueness extends over the whole
real interval $z<1$, $z\ne0$: here
$x=(1-\sqrt{1-z})/2<1/2$, the Euler integral is positive, and the
derivative coefficient is nonzero. The invertible transport therefore
excludes every additional slope. The continuation used to evaluate
two fibers does not alter the ordinary convergence already proved.
\end{proof}

These five classes are distinct from the 35 linear Euler companions.
They agree with \cite[Tables~13 and 15]{CGG}; for the last row the
coefficient convention is
\[
 U_4(n)=(-1)^n\sum_{k=0}^n
   \binom{2k}{k}^{\!2}\binom{2n-2k}{n-k}^{\!2},
\]
which explains the negative denominator. Removing the CM premise
requires quadratic-period independence at the additional fibers
$z=R_\ell c/(c+K_\ell)^2$, $c\in\Z$, $|c|>K_\ell$.
The restricted standard hypothesis does not contain that assertion.

\subsection{The level-nine pullback}\label{sec:nine}
Let
\[
 D(t)=(1-9t)(1+3t),\qquad
 H(t)=D(t)^{-1/2}F_{1/6}\!\left(\frac{1728t^3}{D(t)^3}\right)
      =\sum_{n\ge0}h_nt^n,
\]
where the square root is one at zero. This is
\cite[Theorem~9.2, $\ell=1$]{CCY}, with the factor $1728$ converting
its factorial variable to the present hypergeometric variable.

\begin{theorem}\label{thm:nine}
For a nonzero pair $(a,b)$ and integer $C\ne0$, the weighted series
$\sum h_n(a+bn)C^{-n}$ converges exactly when $|C|\ge20$; its
convergence is then absolute. At each such $C$ there is at most one
algebraic projective slope with value in $\Qbar/\pi$, and a nonzero
pair cannot have value zero. The complete primitive CM list, with
$b>0$, is
\[
\begin{array}{rrr|l}
 a&b&C&\alpha\\\hline
 1&4&27&9\sqrt3/2\\
 7&22&-27&9\sqrt3.
\end{array}
\]
It is exhaustive without a CM premise if quadratic-period independence
holds for the selected non-CM fibers at
$z=1728t^3/D(t)^3$, $t=1/C$, $|C|\ge20$.
\end{theorem}
\begin{proof}
The radius is $\rho=(9+6\sqrt3)^{-1}$. In fact $z(\rho)=1$;
the nonnegative coefficient expansions give $|z(t)|\le z(|t|)<1$
for $|t|<\rho$, whereas $H'(t)$ diverges as $t\uparrow\rho$.
A nonzero polynomial weight preserves the radius, and no integer
reciprocal has modulus $\rho$. This proves convergence.

The transport \eqref{rev:eq:jet-pullback} has
\[
 r=D^{-1/2},\qquad
 u=\frac{3t(1+9t)}D,\qquad v=\frac{3(1+27t^2)}D.
\]
For real $0<|t|<\rho$ the derivative $v$ is nonzero and
$z\in(-1,1)\setminus\{0\}$. The unconditional period theorem gives
uniqueness and nonvanishing, and the stated additional independence
hypothesis gives CM.

At CM, integrality of
\[
 j=\frac{D(t)^3}{t^3}=\left(C-6-\frac{27}C\right)^3
\]
forces its rational cube root to be integral; hence $C\mid27$.
Convergence leaves $C=\pm27$, with $j=8000,-32768$. At these two
parameters, respectively, the transported coefficient pair and scalar
are
\[
 ((a+bu,bv),r)
 =\left(\tfrac35(3,28),\tfrac{3\sqrt{15}}{10}\right),\qquad
   \left(\tfrac38(15,154),\tfrac{3\sqrt6}{8}\right).
\]
The two standard evaluations prove the displayed identities.
\end{proof}

The extra hypothesis cannot be suppressed: before CM is known,
$j=(C-6-27/C)^3$ is generally rational rather than integral.

\subsection{Domb numbers}\label{sec:domb}
For
\[
 D_n=\sum_{k=0}^n\binom nk^2\binom{2k}k\binom{2n-2k}{n-k},
 \qquad H(t)=\sum_{n\ge0}D_nt^n,
\]
Rogers's transformation \cite{RogersDomb} is the germ identity
\begin{equation}\label{eq:domb-pullback}
 H(t)=\frac1{1-4t}F_{1/3}\!\left(\frac{108t^2}{(1-4t)^3}\right).
\end{equation}
Its equivalent coefficient formula is
\[
 D_n=\sum_{k=0}^{\lfloor n/2\rfloor}
   \binom{2k}k^2\binom{3k}k\binom{n+k}{3k}4^{n-2k}.
\]
The singular expansion at $1/16$ gives
$D_n\sim2\,16^n/(\pi^{3/2}n^{3/2})$, with a full expansion in
descending powers of $n$. Thus a nonzero linear weight has absolute
convergence for $|C|>16$, divergence for $|C|<16$, convergence at
$C=16$ exactly for $B=0$, and conditional convergence at $C=-16$
when $B\ne0$. The expansion makes $nD_n/16^n$ eventually decreasing,
so the alternating endpoint is an ordinary sum.

Writing $Y=(C-4)^3/C$, transport gives
\begin{equation}\label{eq:domb-weights}
 (A+B\delta)H(1/C)=\frac{C}{C-4}
 \left(A+\frac{4B}{C-4}
       +\frac{2(C+2)B}{C-4}\thetaop\right)F_{1/3}(108/Y).
\end{equation}
For $|C|\ge16$, $C\ne16$, this map is invertible and $0<108/Y<1$.
The period theorem, with Abel's theorem at $C=-16$, gives unconditional
projective uniqueness. At $C=16$, the only convergent weight is constant
and $H(1/16)=\tfrac43F_{1/3}(1)$, so the endpoint theorem applies.

\begin{proposition}\label{prop:domb-CM}
The primitive Domb identities with integer denominator and a CM fiber
are, up to simultaneous sign,
\[
\begin{array}{rrr|r|l}
 A&B&C&D&\alpha\\\hline
 1&3&-32&-48&2\\
 1&5&64&-60&8\sqrt3/3.
\end{array}
\]
Here $D$ is the endomorphism-order discriminant. Both sums converge
absolutely.
\end{proposition}
\begin{proof}
The two invariants above the level-three parameter $Y$ are
$3$-isogenous; hence at CM their trace is an integer. That trace is
the monic polynomial $Y^3-126Y^2+2944Y$. Since $Y\in\Q$, it follows
that $Y\in\Z$, and
\[
 Y=C^2-12C+48-\frac{64}C
 \quad\Longrightarrow\quad C\mid64.
\]
The convergence filter leaves $C\in\{\pm16,\pm32,\pm64\}$.
Their images form the set
\[
 \{Y\}=\{108,500,686,1458,3375,4913\};
\]
the exact level-three CM test retains $108,1458,3375$, of discriminants
$-12,-48,-60$. The first is the excluded positive endpoint.
Transporting the established standard pairs $(2,15)$ and $(4,33)$
at the other two parameters gives the stated values. Uniqueness
excludes any other coefficient pair.
\end{proof}

The list is unconditional within the CM locus. Exhaustion of all
Domb identities would additionally require quadratic-period independence
at $z=108C/(C-4)^3$, for integer $|C|\ge16$, $C\ne16$. Integrality of
the intermediate denominator was a consequence of CM, and cannot be
used to remove this extra hypothesis.

\section{Modular coordinates and returning correspondences}
\label{sec:moonshine}
For $q=e^{2\pi i\tau}$, put
\[
 (k_\ell,d_\ell)=(24,64),(12,27),(8,16),\qquad
 t_\ell=\left(\frac{\eta(\tau)}{\eta(\ell\tau)}\right)^{k_\ell}
 \quad(\ell=2,3,4).
\]
The Hauptmoduls of \cite[Table~2]{CCY} express the denominator as the
Fricke norm coordinate
\[
 C_1=j,\qquad C_\ell=t_\ell+2d_\ell+d_\ell^2/t_\ell.
\]
Indeed Fricke sends $t_\ell$ to $d_\ell^2/t_\ell$. In the normalized
McKay--Thompson convention $T_g=q^{-1}+O(q)$ \cite{CN},
\begin{equation}\label{eq:moonshine-shifts}
 C_1=T_{1A}+744,\quad C_2=T_{2A}+104,\quad
 C_3=T_{3A}+42,\quad C_4=T_{4A}+24.
\end{equation}
This follows from $t_\ell=q^{-1}-k_\ell+O(q)$. At level four the
elliptic model has invariant $j(2\tau)=(C_4-16)^3/C_4$; the rescaling
of $\tau$ preserves CM.

\begin{proposition}[Integral coordinates and CM]
Every integer $C>64$ with $C\nmid4096$ occurs on the positive convergent
level-four branch and corresponds to a non-CM point.
\end{proposition}
\begin{proof}
The positive continuous function $t_4(iy)$, $y\ge1/2$, takes the
value $16$ at $y=1/2$ and tends to infinity. It therefore takes
$t=(C-32+\sqrt{C(C-64)})/2>16$, at which $C_4=C$. But
\[
 j(2\tau)=C^2-48C+768-4096/C
\]
is rational and nonintegral, whereas a singular modulus is an algebraic
integer.
\end{proof}
For instance, $C=65$ gives $T_{4A}=41$ and $j(2\tau)=117649/65$.
These points satisfy the functional Moonshine identities; no
inverse-$\pi$ evaluation at them is asserted.

\begin{proposition}[A returning correspondence]
Let $C_\ell$ have invariance group $\Gamma_\ell$. If
\[
 C_\ell(g\tau)=C_\ell(\tau),\qquad g\in\mathrm{GL}_2^+(\Q),
\]
and the projective action of $g$ does not belong to $\Gamma_\ell$,
then $\tau$ is imaginary quadratic.
\end{proposition}
\begin{proof}
A Hauptmodul separates points of its quotient, so
$g\tau=\gamma\tau$ for some $\gamma\in\Gamma_\ell$. These groups
have rational projective representatives. The nonscalar rational
matrix $\gamma^{-1}g$ fixes a nonreal point; its fixed-point equation
is therefore a rational quadratic with negative discriminant.
\end{proof}
Replication relates Hecke translates; it does not supply the returning
equality. Deriving such an equality from the scalar period relation
would be a genuine converse theorem.

For precision, the corresponding differential problem has an exact
form. Choose the lift $\tau(z)$ continued from the cusp, with
$z=R_\ell/C_\ell(\tau)\sim R_\ell q$. Set $D=q\,d/dq$ and
$r=\sqrt{1-R_\ell/C}$, continued from $r=1$. The modular
parameterization \cite[Theorem~2.1]{CCY} gives
\[
 F_{s_\ell}(R_\ell/C)=-\frac{DC}{Cr},\qquad
 \frac{\thetaop F}{F}
 =1+\frac{R_\ell}{2(C-R_\ell)}-\frac{CD^2C}{(DC)^2}.
\]
Consequently, at an ordinary point with $DC\ne0$, the original identity
is equivalent to
\begin{equation}\label{eq:moonshine-jet}
 -\frac{C'}{2iCr}
 \left[A+B\left(1+\frac{R_\ell}{2(C-R_\ell)}
                   -\frac{CC''}{(C')^2}\right)\right]=\alpha,
\end{equation}
where primes mean $d/d\tau$. The chosen lift is part of this statement.
Differentiated replication does not by itself impose an additional
relation at the original point.

\section{Arithmetic correspondences and compact Shimura curves}
\label{sec:extensions}
\subsection{The finite-map principle}

The arithmetic part of the argument belongs to finite morphisms of
curves, rather than to any particular hypergeometric equation.

\begin{proposition}[Arithmetic correspondence criterion]
\label{prop:arithmetic-correspondence-packet}\label{prop:extension}
Let $k$ be a number field and let a fixed curve carry an elliptic
invariant $j$ and a parameter $t$ satisfying $P(t,j)=0$, with
$P\in k[T,J]$ of $J$-degree at most $m$. Exclude a specified finite
exceptional set, require $P(t,J)\ne0$ at every allowed parameter, and
assume each fixed $j$ has only finitely many allowed parameters.
If $[k(t):k]\le d$ and the fiber has CM by the order of discriminant
$D$, then
\[
 h(D)\le[k:\Q]md.
\]
There are finitely many such CM parameters. If every admitted identity
forces CM and has at most one rational coefficient line at each
parameter, there are finitely many primitive integer pairs up to sign,
apart from the separately specified exceptional contributions.
\end{proposition}
\begin{proof}
CM reciprocity gives $h(D)=[\Q(j):\Q]$. The tower of residue fields
gives
\[
 [\Q(j):\Q]\le[k(t,j):\Q]
 \le[k:\Q]\,[k(t):k]\,[k(t,j):k(t)]\le[k:\Q]dm.
\]
Siegel's lower bound implies that only finitely many imaginary
quadratic fields have bounded class number. For an order of conductor
$f$ in a field of discriminant $d_0$, the conductor formula gives
\[
 h(f^2d_0)\ge\frac{h(d_0)}3\varphi(f)
             \ge\frac{h(d_0)}3\sqrt{f/2}.
\]
Thus bounded order class number also bounds $f$. There are finitely
many singular moduli, hence finitely many parameters by the finite-fiber
hypothesis. A rational projective line contains exactly one primitive
integer pair up to sign.
\end{proof}

For a modular curve with an honest $\Q$-defined $j$-map, a rational
point has rational $j$, independently of the map's degree. A quotient
requires the degree of the forgotten data. The diagram
\[
\begin{tikzcd}[column sep=large,row sep=large]
 X_0(N) \arrow[r,"j"] \arrow[d,"\pi"'] & \mathbf P^1_j\\
 X_0(N)/W \arrow[r,"t"'] & \mathbf P^1_t
\end{tikzcd}
\]
encodes a correspondence: the norm
$\operatorname{Nm}_{\Q(X_0(N))/\Q(X_0(N)/W)}(J-j)$ is monic of
degree at most $|W|$ in $J$. At rational points of the quotient this
gives $h(D)\le|W|$; if only the value of a degree-$\delta$ coordinate
$t$ is rational, the bound is $|W|\delta$. A Fricke quotient therefore
gives two, whereas larger Atkin--Lehner quotients need not. The maps
must be defined over the asserted number field. Integrality of a
particular coordinate additionally requires its monic arithmetic
equations and is not preserved by every fractional linear change.

\subsection{A class-number-four identity}
\begin{proposition}\label{prop:level6-h4}
Let
\[
 u_6(n)=\binom{2n}{n}\sum_{k=0}^n\binom nk^3,
 \qquad \tau_0=\frac{-3+\sqrt{-21}}6.
\]
Then
\begin{equation}\label{eq:level6-h4}
 \sum_{n\ge0}(2+9n)\frac{u_6(n)}{(-112)^n}
          =\frac{2\sqrt7}{\pi}
\end{equation}
is an absolutely convergent primitive identity. The elliptic parameter
has endomorphism discriminant $-84$ and class number four; its rational
coordinate lies on $X_0(6)/\langle w_2,w_3\rangle$.
\end{proposition}
\begin{proof}
The exact modular inputs in \cite[Tables~2, 4, 15 and 55]{CCY} are
\[
 X_6=\frac{w}{(1-w)^2},\quad
 w=\left(\frac{\eta(\tau)\eta(6\tau)}
                   {\eta(2\tau)\eta(3\tau)}\right)^{12},\quad
 G_6(X)=(1+4X)(1-32X).
\]
The associated weight-two form has coefficients $u_6(n)$. At
$q=e^{2\pi i\tau_0}=-e^{-2\pi\sqrt{7/12}}$, the tabulated exact
values are $X_6=-1/112$ and $\lambda_6=2/9$. Thus their
Theorem~2.4, equation~(2.14), yields
\[
 \sum_{n\ge0}(2+9n)u_6(n)X_6^n
 =\frac9{2\pi\sqrt{G_6(X_6)}}\sqrt{\frac{12}7}
 =\frac{2\sqrt7}{\pi},
\]
since $\sqrt{G_6(-1/112)}=9\sqrt3/14$.
The bound $u_6(n)\le4^n(\sum_k\binom nk)^3=32^n$ proves absolute
convergence, and $\gcd(2,9)=1$.

The primitive equation $6\tau_0^2+6\tau_0+5=0$ gives discriminant
$-84$. Its reduced forms are
\[
 [1,0,21],\quad[2,2,11],\quad[3,0,7],\quad[5,4,5];
\]
the bound
$a\le\sqrt{84/3}<6$ proves completeness. Hence its order
$\mathcal O=\Z[\sqrt{-21}]$ has four classes.

The quotient explains the rational coordinate intrinsically. Put
$s=\sqrt{-21}$ and let
\[
 \mathfrak p_2=(2,1+s),\quad \mathfrak p_3=(3,s),\quad
 \mathfrak n=\mathfrak p_2\mathfrak p_3=(6,3+s).
\]
The two ramified ideal classes have order two and are independent:
the norm $a^2+21b^2$ represents none of $2,3,6$. They therefore
generate the class group. On $\C/\mathfrak a$, the unique
endomorphism-stable subgroup of order six is
$\mathfrak n^{-1}\mathfrak a/\mathfrak a$; quotienting by its
ramified $p$-part realizes $w_p$ as multiplication by
$[\mathfrak p_p^{-1}]$. Hence all four CM pairs have the same image
on the quotient. Their set is Galois stable, so this image is rational.
For the standard subgroup at $\tau_0$, stability follows from
$s=6\tau_0+3$, which preserves $\langle1/6\rangle$ modulo the period
lattice.
\end{proof}
This example makes the quotient-degree bound necessary; it gives no
converse for arbitrary level-six identities.

\subsection{Local monodromy and compact triangle systems}
The compact Shimura curve $X_6^*$ has triangle signature $(4,2,6)$.
Its Clausen identity is
\[
 {}_2F_1(1/24,5/24;3/4;z)^2
 ={}_3F_2(1/12,1/4,5/12;1/2,3/4;z).
\]
The rank-two equation has finite local projective monodromy of orders
$4,2,6$, and infinite global hyperbolic triangle monodromy. This
distinguishes its realization from the elliptic families above.

\begin{proposition}[The compact monodromy obstruction]
\label{prop:compact-monodromy-obstruction}
Let $L$ be a rank-two complex local system on a smooth connected
algebraic curve, with finite projective local monodromy at every point
of the boundary in its smooth compactification. After any finite
algebraic base change, no positive symmetric power of $L$ is
projectively equivalent, on a nonempty algebraic open set, to the same
symmetric power of $R^1f_*\C$ for a non-isotrivial elliptic family.
Rank-one twists and algebraic meromorphic gauges do not alter this
conclusion.
\end{proposition}
\begin{proof}
Pass to a common finite cover and compactify. The nonconstant elliptic
invariant $j$ has a pole. After a finite local base change the elliptic
family has multiplicative semistable reduction there, with monodromy
$\bigl(\begin{smallmatrix}1&n\\0&1\end{smallmatrix}\bigr)$,
$n\ne0$. Every positive symmetric power is a nonidentity unipotent
matrix, hence has infinite projective order in characteristic zero.
The pulled-back $L$ has finite projective local monodromy at that
point: it is a power of an original boundary monodromy, or trivial
over an interior point. This is a contradiction. Twists disappear
projectively, and gauges conjugate monodromy.
\end{proof}

An isotrivial elliptic family also cannot realize this compact triangle
system: its projective monodromy is finite. For comparison, the standard
Gauss equation has a repeated exponent zero at the cusp and Wronskian
$cz^{-1}(1-z)^{-1/2}$, $c\ne0$, so it has nontrivial unipotent
monodromy there. Every nonconstant finite cover retains such a
monodromy above the cusp. Thus the compact and standard systems, or
their equal positive symmetric powers, cannot become projectively
equivalent on a common finite algebraic cover. The obstruction is
global; it neither excludes an isolated special-value identity nor a
quaternionic abelian-surface realization.

\subsection{CM degrees without an elliptic invariant}
\label{subsec:compact-class-field}
\begin{proposition}[Compact CM degree bound]
\label{prop:compact-cm-degree}
Let $D>1$ be the reduced discriminant of an indefinite quaternion
algebra over $\Q$, and let $N$ be square-free and coprime to $D$.
On the canonical model $X=X_0(D,N)$, let $W$ be an Atkin--Lehner
subgroup, $Y=X/W$, and let $t:Y\to\mathbf P^1$ be a nonconstant
$\Q$-morphism of degree $\delta$. If $P\in X(\Qbar)$ is CM by
an optimally embedded order $R$, and $Q=\pi(P)$, then
\begin{equation}\label{eq:compact-class-bound}
 h(R)\le |W|\delta[\Q(t(Q)):\Q].
\end{equation}
Bounded parameter degree therefore gives finitely many CM points and
parameters on this fixed curve.
\end{proposition}
\begin{proof}
Put $K=\operatorname{Frac}(R)$ and let $H_R$ be its ring class field.
For the stated canonical models and optimal embeddings,
Gonz\'alez--Rotger \cite[Theorem~5.8(1)]{GonzalezRotger} prove
\begin{equation}\label{eq:GR-class-field}
 H_R=K\Q(P),\qquad [H_R:K]=h(R).
\end{equation}
Their reciprocity identifies this Galois action with the invertible
ideal-class action, free on each fixed local-orientation branch by
their Proposition~5.6. It is a statement about the field of the
individual coarse point, including elliptic stabilizers.
The field diagram
\[
\begin{tikzcd}[column sep=large,row sep=large]
 \Q(P) \arrow[r,hook] & H_R=K\Q(P)\\
 \Q(Q) \arrow[u,hook] & K \arrow[u,hook]\\
 \Q(t(Q)) \arrow[u,hook] & \Q \arrow[u,hook] \arrow[ul,hook]
\end{tikzcd}
\]
and the two finite maps give
\[
 h(R)\le[\Q(P):\Q]
       \le|W|[\Q(Q):\Q]
       \le|W|\delta[\Q(t(Q)):\Q].
\]
There is no extra factor two from adjoining $K$. Bounded class number
gives finitely many orders, and each order has finitely many optimal
embedding classes by the same Proposition~5.6.
\end{proof}

For $X_6^*=X_0(6,1)/W_6$, one has $|W_6|=4$. The coordinate taking
$0,1,\infty$ at its unique CM points of discriminants $-4,-24,-3$
is defined over $\Q$: these singleton CM sets are Galois stable.
It has degree one, so $h(R)\le4[\Q(t):\Q]$. Equality occurs at a
rational parameter. Indeed $R=\Z[\sqrt{-30}]$ has discriminant
$-120$ and reduced forms
\[
 [1,0,30],\quad[2,0,15],\quad[3,0,10],\quad[5,0,6].
\]
The ramified classes above $2$ and $3$ independently generate its
four-element class group, since $a^2+30b^2$ represents none of
$2,3,6$. Their ideal actions are the Atkin--Lehner actions by
\cite[Lemma~5.9]{GonzalezRotger}. Here the local-orientation factors
are trivial, so these four points form the entire CM set. Its quotient
is therefore a Galois-stable singleton. This is the $-120$ row of
\cite[Theorem~1]{Yang}.

The optimal-embedding hypothesis concerns the order itself. For example,
$\Q(\sqrt{-3})$ embeds in the discriminant-six algebra, but its order
of conductor two does not optimally embed in a maximal order: at a
ramified quaternion prime the intersection with an embedded quadratic
field is its maximal local order. No assertion for arbitrary levels,
orders, or totally real ground fields is implicit here.

\subsection{The period target belongs to the realization}
\label{subsec:compact-period-target}
Suppose a specified rank-two comparison matrix is
\[
 P=\begin{pmatrix}\omega&\omega_2\\\eta&\eta_2\end{pmatrix},
 \qquad \det P=c(2\pi i),\qquad c\in\Qbar^\times.
\]
For a quaternionic abelian surface this must be a specified multiplicity
summand after splitting the quaternion action over $\Qbar$, together
with its algebraic polarization. A differential equation alone does
not provide this realization.

\begin{proposition}[Conditional rank-two mechanism]
\label{prop:rank-two-torsor-cm}
Assume that at every non-CM fiber the numerical comparison point is
Zariski generic over $\Qbar$ in its full $\mathrm{GL}_2$ comparison
torsor, with the determinant giving the Tate comparison coordinate.
Then
\[
 a\omega^2+b\omega\eta=\alpha\pi,
 \qquad a,b,\alpha\in\Qbar,\quad(a,b)\ne(0,0),
\]
forces CM.
\end{proposition}
\begin{proof}
Trivialize the torsor over $\Qbar$. The left side is a nonzero quadratic
function of the first column; the right side is
$\alpha\det(P)/(2ic)$. These cannot agree identically on
$\mathrm{GL}_2$: if $\alpha\ne0$, fix the first column and vary the
second; if $\alpha=0$, vary the first column. Genericity excludes
their equality at the comparison point.
\end{proof}
The determinant varies on the torsor; it is not a fixed numerical
$2\pi i$ defining an algebraic variety over $\Qbar$. Genericity is
the extra period-conjecture input, beyond the computation of the
Mumford--Tate group or of its tensor invariants.

Now suppose the chosen branch has a verified normalization
\begin{equation}\label{eq:compact-normalized-square}
 \mathcal F(t)=\frac{r(t)\omega(t)^2}{\Omega_0^2},\qquad
 \frac{\thetaop\mathcal F}{\mathcal F}
       =u(t)+v(t)\frac{\eta(t)}{\omega(t)},\qquad rv\ne0,
\end{equation}
with algebraic $r,u,v$ at the parameter. Transport identifies its
coefficient space with $a\omega^2+b\omega\eta$, and the natural
target is
\begin{equation}\label{eq:compact-natural-target}
 (A+B\thetaop)\mathcal F(t)\in\Qbar\,\frac{\pi}{\Omega_0^2}.
\end{equation}
Conditional CM forcing and the compact degree bound give finitely many
bounded-degree parameters on a fixed compact curve. For finiteness of
primitive pairs one must also prove slope uniqueness. It follows if
$\eta/\omega$ is transcendental: two independent pairs eliminate
$\pi$ and give an algebraic linear relation between $\omega$ and
$\eta$. Chudnovsky supplies this when the CM realization reduces
algebraically to elliptic de Rham cohomology, with $\omega$ a nonzero
holomorphic period on a rational homology cycle and $\eta$ an independent
de Rham class evaluated on that same cycle. An arbitrary algebraic
combination of CM cycles may instead be an endomorphism eigencycle;
the rational-cycle condition must be retained. Singular parameters
and vanishing derivative coefficients require separate treatment.

\begin{proposition}[Weight obstruction]
\label{prop:compact-weight-obstruction}
Assume $\omega,\eta,\Omega_0$ are actual degree-one periods of abelian
varieties over $\Qbar$, and assume genericity of their numerical
comparison point in the tensor torsor generated by these varieties
and the Tate object. Then no value of
\eqref{eq:compact-normalized-square} with a linear weight can equal
$\alpha/\pi$ for $\alpha\in\Qbar^\times$.
\end{proposition}
\begin{proof}
Such an equality says
$\pi(a\omega^2+b\omega\eta)=\alpha\Omega_0^2$.
The central weight homotheties multiply degree-one periods by
$\lambda$ and the Tate period by $\lambda^2$. The two sides have
weights four and two. Genericity would turn the equality into an
identity on the torsor, incompatible with these distinct characters
and the nonzero right side.
\end{proof}
For the standard cusp normalization, the denominator is $\pi^2$,
of weight four, so this obstruction does not apply. For a compact
normalization by an actual degree-one period squared, both
$\mathcal F$ and its natural target $\pi/\Omega_0^2$ have weight zero.

This difference is already visible unconditionally in Yang's example.
Let
\[
 d_n=\frac{(1/12)_n(1/4)_n(5/12)_n}{(1/2)_n(3/4)_n n!},
 \qquad \Omega_{-4}=\sqrt\pi\,\frac{\Gamma(1/4)}{\Gamma(3/4)}.
\]
\begin{proposition}[A compact CM value outside $\Qbar/\pi$]
The primitive, absolutely convergent identity
\begin{equation}\label{eq:yang-example}
 \sum_{n\ge0}(9+598n)d_n\left(\frac{8748}{15625}\right)^n
     =90\sqrt5\,\frac{\pi}{\Omega_{-4}^2}
\end{equation}
has right side outside $\Qbar/\pi$.
\end{proposition}
\begin{proof}
The $D=-52$ row of \cite[Theorem~1]{Yang} gives
$(M,N,R_1,R_2,R_3)=(2^2 3^7,5^6,64584,972,1)$ and multiplier
$4M^{3/4}N^{1/4}\pi/(12^{1/4}\Omega_{-4}^2)$.
Dividing the coefficients and multiplier by $108$ gives
\eqref{eq:yang-example}. Its parameter is in $(0,1)$.
The beta integral identifies the normalization as an actual CM period:
\[
 \Omega_{-4}=B(1/4,1/2)
       =2\int_1^\infty\frac{dX}{\sqrt{X^3-X}}.
\]
Chudnovsky's theorem, in \cite[Corollary~6]{Waldschmidt}, makes this
nonzero period of $y^2=X^3-X$ algebraically independent of $\pi$.
Membership of the right side in $\Qbar/\pi$ would contradict that
independence.
\end{proof}

\subsection{The exact role of the construction theorems}
\label{subsec:triangle-source-packets}
The external constructions enter the theory at distinct stages.
Chen--Glebov--Goenka \cite[Theorems~4.1--4.2 and 6.1--6.9]{CGG}
identify the elliptic Picard--Fuchs equations and their hypergeometric
periods on specified simply connected domains. Their specialization
theorems require a quadratic parameter
$\tau=(-b+\sqrt{-d})/(2a)$, as stated in their equation~(22).
For such $\tau$ with $E_6(\tau)\ne0$, their Theorem~8.1 proves
\[
 \frac{E_4(\tau)}{E_6(\tau)}
 \left(E_2(\tau)-\frac3{\pi\operatorname{Im}\tau}\right)
       \in\Q(j(\tau)),
\]
and supplies an integer clearing the denominator. At $j=1728$ this
quotient is undefined; their precursor identity~(21) is used instead.

The arithmetic mechanism is a torsion trace. On an algebraic model,
for the nonreal endomorphism $\mu=a\tau$ of its period lattice,
their equations~(28)--(33) give
\[
 \mu(\mu-\bar\mu)e_2^*(\Lambda)
       =\sum_{0\ne P\in E[(\mu)]}x(P).
\]
The trace is algebraic, the nonzero multiplier permits solving for
$e_2^*$, and conjugation descends the coefficient to $\Q(j)$.
On a non-CM curve every multiplier is an integer, so this argument
loses precisely the factor $\mu-\bar\mu$. It computes the CM
coefficient; it does not prove the converse.

Yang's compact theorem instead supplies the first-solution constant
\[
 C_1=\frac4{12^{1/4}}\frac{\pi}{\Omega_{-4}^2},
\]
and $C_1^{-1}$ for the second local solution. Its Hecke and algebraic
equations evaluate derivative coefficients at specified CM points.
The reciprocal constants make the local branch part of the problem,
and \eqref{eq:yang-example} prevents replacing this target by
$\Qbar/\pi$ automatically.

Takeuchi's classification \cite{Takeuchi} classifies arithmetic triangle
groups, rather than the CM points at rational uniformizer values.
Similarly, \cite[Theorem~2.2]{AG} uses three compatible equations
on Frobenius solutions $w_0=y_0^2$, $w_1=y_0y_1$, $w_2=y_1^2/2$:
\[
 aw_0+b\thetaop w_0=1/\pi,\qquad
 aw_1+b\thetaop w_1=ir,\qquad
 aw_2+b\thetaop w_2=-\pi(\beta+r^2)/2.
\]
Here $\beta$ denotes their auxiliary parameter, distinct from the
multiplier of our original identity. The first equation does not
imply the other two. Their $1/\pi^2$ Calabi--Yau construction uses
a different system and explicitly conjectural specialization
properties after their equation~(41). In each case the valid
sequence is realization, a separate CM criterion, reciprocity, and
CM evaluation; no construction theorem removes the missing scalar
CM implication.
